\section*{अध्याय \ref{ch:g10:reaction-progress-table} --- अभिक्रिया की प्रगति सारणी}

\begin{solution}{exo:g10:reaction-progress-table:1}
\ce{H2}: $4.0 - 2x$; \ce{O2}: $3.0 - x$; \ce{H2O}: $2x$। डाइहाइड्रोजन
$x = 2.0$ पर ख़त्म होती है, डाइऑक्सीजन $x = 3.0$ पर: $x_{\max} = \qty{2.0}{mol}$,
डाइहाइड्रोजन सीमांत है। \omterm{def:g10:reaction-progress-table:system}{अंतिम अवस्था}: \qty{0}{mol} \ce{H2},
\qty{1.0}{mol} \ce{O2}, \qty{4.0}{mol} \ce{H2O}।
\end{solution}

\begin{solution}{exo:g10:reaction-progress-table:2}
कार्बन $x = 3.0$ पर ख़त्म होगा, डाइऑक्सीजन $x = 2.0$ पर:
$x_{\max} = \qty{2.0}{mol}$, डाइऑक्सीजन सीमांत है। \omterm{def:g10:reaction-progress-table:system}{अंतिम अवस्था}:
\qty{1.0}{mol} कार्बन, कोई डाइऑक्सीजन नहीं, \qty{2.0}{mol} कार्बन डाइऑक्साइड।
\end{solution}

\begin{solution}{exo:g10:reaction-progress-table:3}
$x_{\max} = \qty{0.30}{mol}$ (सल्फ़र सीमांत)। \omterm{def:g10:reaction-progress-table:system}{अंतिम अवस्था}:
\qty{0.20}{mol} लोहा, कोई सल्फ़र नहीं, \qty{0.30}{mol} आयरन सल्फ़ाइड।
\end{solution}

\begin{solution}{exo:g10:reaction-progress-table:4}
डाइनाइट्रोजन $x = 1.0$ पर ख़त्म होती है, डाइहाइड्रोजन $x = 2.4/3 = 0.80$ पर:
$x_{\max} = \qty{0.80}{mol}$। \omterm{def:g10:reaction-progress-table:system}{अंतिम अवस्था}: \qty{0.20}{mol} \ce{N2}, कोई
\ce{H2} नहीं, $2 \times 0.80 = \qty{1.6}{mol}$ \ce{NH3}।
\end{solution}

\begin{solution}{exo:g10:reaction-progress-table:5}
(ख): $4/2 = 2/1$। (क) में कार्बन मोनोऑक्साइड सीमांत है, (ग) में भी।
\end{solution}

\begin{solution}{exo:g10:reaction-progress-table:6}
$n(\ce{CH4}) = 32 / 16.0 = \qty{2.0}{mol}$, इसलिए \omterm{def:g10:reaction-progress-table:stoichiometric}{रससमीकरणमितीय मिश्रण} के लिए
$x_{\max} = \qty{2.0}{mol}$। डाइऑक्सीजन: $2 \times 2.0 = \qty{4.0}{mol}$,
यानी $4.0 \times 32.0 = \qty{128}{g}$। पानी:
$\qty{4.0}{mol} \times \qty{18.0}{g/mol} = \qty{72}{g}$।
\end{solution}

\begin{solution}{exo:g10:reaction-progress-table:7}
$n(\ce{C3H8}) = 1000 / 44.0 = \qty{22.7}{mol}$; कार्बन डाइऑक्साइड
$3 \times 22.7 = \qty{68.2}{mol}$; $V = 68.2 \times 24.1 =
\qty{1.64e3}{L}$, लगभग \qty{1.6}{m^3}।
\end{solution}

\begin{solution}{exo:g10:reaction-progress-table:8}
$x = \qty{1.0}{mol}$ पर: \qty{1.0}{mol} \ce{CH4}, \qty{1.0}{mol}
\ce{O2}, \qty{1.0}{mol} \ce{CO2}, \qty{2.0}{mol} \ce{H2O}। डाइऑक्सीजन
$x = \qty{1.5}{mol}$ पर खप जाती है: अभिक्रिया वहीं रुक जाती है, इसलिए इससे बड़ी
प्रगति तक कभी नहीं पहुँचा जाता।
\end{solution}

\begin{solution}{exo:g10:reaction-progress-table:9}
$n_0(\ce{Zn}) = 1.30 / 65.4 = \qty{0.0199}{mol}$;
$n_0(\ce{Cu^{2+}}) = 0.10 \times 0.1000 = \qty{0.0100}{mol}$। कॉपर
\omterm{def:g9:ions:ion}{आयन} सीमांत हैं: $x_{\max} = \qty{0.0100}{mol}$। बना हुआ ताँबा:
$0.0100 \times 63.5 = \qty{0.635}{g}$; बचा हुआ ज़िंक:
$(0.0199 - 0.0100) \times 65.4 = \qty{0.65}{g}$। कोई कॉपर \omterm{def:g9:ions:ion}{आयन} नहीं बचता:
नीला रंग चला गया है।
\end{solution}

\begin{solution}{exo:g10:reaction-progress-table:10}
$n_0(\ce{Mg}) = 0.48 / 24.3 = \qty{0.0198}{mol}$, जो
$x = 0.0198$ पर ख़त्म होगा; $n_0(\ce{H+}) = \qty{0.020}{mol}$, जो
$x = 0.010$ पर ख़त्म होगा। अम्ल सीमांत है: $x_{\max} = \qty{0.010}{mol}$, इसलिए
\qty{0.010}{mol} डाइहाइड्रोजन, $0.010 \times 24.1 = \qty{0.24}{L}$।
\end{solution}

\begin{solution}{exo:g10:reaction-progress-table:11}
\qty{0.30}{mol} सल्फ़र वाले \omterm{def:g10:reaction-progress-table:stoichiometric}{रससमीकरणमितीय मिश्रण} को \qty{0.30}{mol} लोहा
चाहिए; \qty{0.20}{mol} अधिक है, यानी
$0.20 / 0.30 \approx \qty{67}{\percent}$ आधिक्य।
\end{solution}

\begin{solution}{exo:g10:reaction-progress-table:12}
बराबर मात्राएँ: $m(\ce{Fe})/55.8 = m(\ce{S})/32.1$, और
$m(\ce{Fe}) + m(\ce{S}) = \qty{10.0}{g}$। इसलिए
$m(\ce{Fe}) = 10.0 \times 55.8 / (55.8 + 32.1) = \qty{6.35}{g}$ और
$m(\ce{S}) = \qty{3.65}{g}$।
\end{solution}

\begin{solution}{exo:g10:reaction-progress-table:13}
\begin{enumerate}
  \item $n_0(\ce{CaCO3}) = 10.0 / 100.1 = \qty{0.0999}{mol}$, जो अकेला
        \omterm{def:g7:chemical-reactions:reactant}{अभिकारक} है: $x_{\max} = \qty{0.0999}{mol}$; इससे \ce{CaO} और \ce{CO2}
        का \qty{0.0999}{mol} मिलता है, यानी
        $0.0999 \times 24.1 = \qty{2.41}{L}$ कार्बन डाइऑक्साइड।
  \item $n_0(\ce{CaO}) = \qty{0.0999}{mol}$, $n_0(\ce{H2O}) = 1.8 / 18.0
        = \qty{0.100}{mol}$: बिना बुझा चूना (बस थोड़ा-सा) सीमांत है;
        \omterm{def:g6:pure-substances-and-mixtures:mixture}{मिश्रण} लगभग रससमीकरणमितीय है। \ce{Ca(OH)2}:
        $M = 40.1 + 2 \times (16.0 + 1.0) = \qty{74.1}{g/mol}$, द्रव्यमान
        $0.0999 \times 74.1 = \qty{7.40}{g}$।
\end{enumerate}
\end{solution}

\begin{solution}{exo:g10:reaction-progress-table:14}
हर फ़्लास्क में $n_0(\ce{H+}) = \qty{0.050}{mol}$; यह
$x = \qty{0.025}{mol}$ पर ख़त्म होगा। मैग्नीशियम: 0.15, 0.30, 0.45, 0.60, 0.75,
\qty{0.90}{g} से 0.0062, 0.0123, 0.0185, 0.0247, 0.0309,
\qty{0.0370}{mol} मिलते हैं। पहले चार फ़्लास्कों में मैग्नीशियम सीमांत है
और डाइहाइड्रोजन $n(\ce{Mg}) \times 24.1$ है: 0.15, 0.30, 0.45 और
\qty{0.60}{L}। आख़िरी दो में अम्ल सीमांत है: $0.025 \times 24.1
= \qty{0.60}{L}$। आलेख लगभग \qty{0.61}{g} मैग्नीशियम तक मूलबिंदु से गुज़रती
एक सीधी रेखा है, फिर \qty{0.60}{L} पर एक क्षैतिज
पठार।
\end{solution}

\begin{solution}{exo:g10:reaction-progress-table:15}
$n_0(\text{एथेनॉल}) = 4.6 / 46.0 = \qty{0.100}{mol}$ ($x = 0.100$ पर
ख़त्म होगा); $n_0(\ce{O2}) = 4.8 / 24.1 = \qty{0.199}{mol}$
($x = 0.199/3 = 0.0664$ पर ख़त्म होगा)। डाइऑक्सीजन सीमांत है,
$x_{\max} = \qty{0.0664}{mol}$। बचा हुआ एथेनॉल: $0.100 - 0.0664 =
\qty{0.0336}{mol}$, यानी $0.0336 \times 46.0 = \qty{1.5}{g}$। कार्बन
डाइऑक्साइड: $2 \times 0.0664 = \qty{0.133}{mol}$, यानी
$0.133 \times 24.1 = \qty{3.2}{L}$।
\end{solution}

\begin{solution}{pb:g10:reaction-progress-table:1}
\textbf{1.} बाईं ओर: 2 Na, 6 N। दाईं ओर: 2 Na, $3 \times 2 = 6$ N। संतुलित।

\textbf{2.} बाईं ओर: 10 Na, 2 K, 2 N, 6 O। दाईं ओर: K: 2; Na: $5 \times 2 =
10$; O: $1 + 5 = 6$; N: 2। संतुलित।

\textbf{3.} \omterm{def:g9:periodic-table-first-look:periodic-table}{आवर्त सारणी} वाला अध्याय (\cref{ch:g9:periodic-table-first-look}):
सोडियम पानी के साथ तीव्र अभिक्रिया करके एक संक्षारक विलयन और
डाइहाइड्रोजन देता है। दुर्घटना के बाद थैली त्वचा, आँखों और साँस की नम \omterm{def:g4:air-a-mixture-of-gases:air}{वायु} के
संपर्क में रहती है।

\textbf{4.} $M(\ce{NaN3}) = 23.0 + 3 \times 14.0 = \qty{65.0}{g/mol}$;
$M(\ce{KNO3}) = 39.1 + 14.0 + 3 \times 16.0 = \qty{101.1}{g/mol}$।

\textbf{5.} \ce{NaN3}: $1.00 - 2x$; \ce{Na}: $2x$; \ce{N2}: $3x$।

\textbf{6.} $1.00 - 2x = 0$, यानी $x_{\max} = \qty{0.500}{mol}$; सोडियम
ऐज़ाइड, जो अकेला \omterm{def:g7:chemical-reactions:reactant}{अभिकारक} है, सीमांत है।

\textbf{7.} \ce{Na}: $2 \times 0.500 = \qty{1.00}{mol}$; \ce{N2}:
$3 \times 0.500 = \qty{1.50}{mol}$।

\textbf{8.} $1.50 \times 24.1 = \qty{36.2}{L}$।

\textbf{9.} \ce{Na}: $1.00 - 10x$; \ce{KNO3}: $n - 2x$; \ce{K2O}: $x$;
\ce{Na2O}: $5x$; \ce{N2}: $x$।

\textbf{10.} $1.00/10 = n/2$, इसलिए $n = \qty{0.200}{mol}$, यानी
$0.200 \times 101.1 = \qty{20.2}{g}$।

\textbf{11.} $x_{\max} = \qty{0.100}{mol}$: कोई सोडियम नहीं, कोई पोटैशियम
नाइट्रेट नहीं, \qty{0.100}{mol} \ce{K2O}, \qty{0.500}{mol} \ce{Na2O},
\qty{0.100}{mol} \ce{N2}।

\textbf{12.} सोडियम $x = 0.100$ पर ख़त्म होगा, पोटैशियम नाइट्रेट
$x = 0.150/2 = 0.075$ पर: नाइट्रेट सीमांत है, $x_{\max} =
\qty{0.075}{mol}$, और $1.00 - 10 \times 0.075 = \qty{0.25}{mol}$ सोडियम
\omterm{def:g9:periodic-table-first-look:metal}{धातु} थैली में बची रहेगी, नमी के साथ अभिक्रिया के लिए तैयार।

\textbf{13.} सोडियम ऐज़ाइड के प्रति \omterm{def:g10:the-mole:amount}{मोल} $1.50 + 0.100 = \qty{1.60}{mol}$
नाइट्रोजन।

\textbf{14.} $n = 60 / 24.1 = \qty{2.49}{mol}$।

\textbf{15.} $2.49 / 1.60 = \qty{1.56}{mol}$ सोडियम ऐज़ाइड।

\textbf{16.} $1.56 \times 65.0 = \qty{101}{g}$।

\textbf{17.} $0.200 \times 1.56 = \qty{0.311}{mol}$, यानी
$0.311 \times 101.1 = \qty{31.5}{g}$ पोटैशियम नाइट्रेट।

\textbf{18.} $2.49 / 1.50 = \qty{1.66}{mol}$, यानी
$1.66 \times 65.0 = \qty{108}{g}$: दूसरी अभिक्रिया लगभग
\qty{7}{g} सोडियम ऐज़ाइड बचाती है।

\textbf{19.} \omterm{def:g10:reaction-progress-table:limiting}{पूर्ण अभिक्रिया} तभी रुकती है जब उसका \omterm{def:g10:reaction-progress-table:limiting}{सीमांत अभिकारक} पूरी तरह खप
जाए; यहाँ सोडियम ऐज़ाइड अकेला \omterm{def:g7:chemical-reactions:reactant}{अभिकारक} है, इसलिए \omterm{def:g10:reaction-progress-table:system}{अंतिम अवस्था} में
उसकी मात्रा $1.00 - 2x_{\max} = 0$ है। थैली में कुछ भी विषैला नहीं बचता।

\textbf{20.} लगभग \qty{101}{g} सोडियम ऐज़ाइड \qty{60}{L} का एयरबैग
भरता है।
\end{solution}
